\section{从 text-to-text 到 omni model：所有模态都要先变成 token}

前面的课程始终围绕语言 token。现在问题变成：图像、视频、音频和其他连续信号怎样进入 Transformer？当前主流路线没有放弃 Transformer，而是先把非文本模态转换为离散或连续 tokens，再与语言 tokens 一起建模。

官方可执行讲义把本讲标为 \texttt{Lecture 17: multimodal models}。这里沿用同一主线，但把模型名称重组为四个可比较问题：representation 怎样得到、alignment 怎样学习、fusion 在哪里发生、generation 由谁完成。

\begin{figure}[H]
\centering
\includegraphics[width=0.78\textwidth]{multimodality.png}
\caption{现实世界是多模态的：理解与生成都需要跨越 text-only interface。}
\figsource{Stanford CS336 2026 Lecture 17 官方可执行讲义，\texttt{images/multimodality.png}。}
\end{figure}

\begin{knowledgebox}{读图：多模态系统有两个不同问题}
\textbf{Understanding} 要把 image/video/audio 编码成模型可消费的表示；\textbf{generation} 要把模型内部表示解码回 pixels、frames 或 waveforms。很多 VLM 只解决前者：能看图回答，但不能原生生成高质量图像。
\end{knowledgebox}

\begin{importantbox}{统一抽象}
多模态模型仍可写成 ``encoder $\rightarrow$ tokens $\rightarrow$ Transformer $\rightarrow$ decoder''。真正的差异在于 token 是连续还是离散、数量有多少、如何保留空间/时间结构，以及训练 loss 怎样平衡不同模态。
\end{importantbox}

\begin{knowledgebox}{Nota de clase}
Percy 把目标称为 omni model：输入和输出都能接受任意模态组合。当前系统大多只实现其中一部分，因此理解模型宣传时要问清楚“支持输入”还是“支持生成”。
\end{knowledgebox}

\subsection{术语表：representation、alignment、fusion 与 generation}

为了让后续模型比较不被品牌名称淹没，本节先固定四个术语。它们分别回答“非文本信号如何表示”“图文为何能对应”“视觉信息如何进入 LM”“模型怎样回到像素或波形”；同一个系统可以在其中三项很强，却仍不具备第四项。

\begin{table}[H]
\centering
\begin{tabular}{L{0.19\textwidth}L{0.31\textwidth}L{0.4\textwidth}}
\toprule
术语 & 核心问题 & 典型实现 \\
\midrule
Representation & 图像/视频怎样变成有限数量 tokens & ViT patches、CLIP features、VQ codebook indices \\
Alignment & 哪些视觉表示与哪些文本语义对应 & CLIP contrastive loss、SigLIP pairwise sigmoid \\
Fusion & 视觉 tokens 在哪里影响语言计算 & input concatenation、cross-attention、DeepStack \\
Generation & 内部状态怎样还原成可见输出 & diffusion decoder、VQ decoder、native image-token LM \\
\bottomrule
\end{tabular}
\caption{多模态模型的四个基本对象。}
\end{table}

\subsection{Token budget 是第一性约束}

本节从容量约束建立直觉。文本一个 token 往往承载一个 subword 的语义，而图像 patch 和视频 frame 的信息密度不同。高分辨率图像会迅速消耗 context window；视频还多一个时间轴。多模态 architecture 的很多设计——patch merging、dynamic resolution、frame sampling——本质上都在控制 token budget。

若 patch size 为 $P$，空间 merge factor 为 $m$，保留 $F$ 帧视频，则视觉 token 数可以粗略写成

$$
N_{\mathrm{image}}\approx
\left\lceil\frac{H}{P}\right\rceil
\left\lceil\frac{W}{P}\right\rceil\frac{1}{m^2},
\qquad
N_{\mathrm{video}}\approx F\,N_{\mathrm{frame}}.
$$

其中 $H,W$ 表示分辨率，$m=2$ 对应 $2\times2$ patch merge。公式的意义不是精确复现某个模型，而是显示两个事实：分辨率翻倍会让二维 token 数近似增至四倍；视频长度又把每帧成本乘上 $F$。因此，视觉细节、帧率和可用于文本推理的剩余 context 之间必然竞争。

\begin{warningbox}{Token budget 也是信息分配}
压缩视觉 tokens 不只是省显存。过度压缩会先损失 OCR、小目标和细粒度空间关系；保留过多视觉 tokens 又会挤占问题、工具轨迹和答案的上下文。评测必须按任务类型观察，而不能只看平均准确率。
\end{warningbox}

\subsection{本章小结}

本节建立了统一视角：先编码，再让 Transformer 处理 tokens。接下来先看 CLIP/SigLIP 怎样学习可对齐的视觉表示。

